\chapter{Efek Elektronik dan Zat Antara Reaktif}\label{ch:b1:electronic-effects}

Asam trifluoroetanoat, \ce{CF3COOH}, sekitar sepuluh ribu kali lebih kuat
daripada asam etanoat, \ce{CH3COOH}, walaupun ikatan O--H yang asam
berjarak tiga ikatan dari atom-atom fluorin. Fenol sejuta kali lebih
asam daripada etanol, dan molekul amonia yang terikat pada cincin
benzena --- anilina --- sejuta kali kurang basa daripada yang terikat
pada gugus metil. Kimia organik penuh dengan kontras semacam itu, dan dua
efek menjelaskan sebagian besar di antaranya: tarikan atom-atom
elektronegatif di sepanjang ikatan $\sigma$, dan tersebarnya
pasangan-pasangan elektron di atas sistem $\pi$. Kedua efek yang sama
menentukan ikatan mana yang putus, zat antara mana yang terbentuk, dan
di mana pereaksi menyerang: keduanya adalah tata bahasa \omterm{def:b1:elementary-steps:elementary-step}{mekanisme
reaksi} pada bab-bab berikutnya.

\begin{recall}
\omterm{def:b1:periodicity:electronegativity}{Keelektronegatifan} dan kecenderungannya (\cref{ch:b1:periodicity});
\omterm{def:b1:lewis-resonance-vsepr:lewis}{struktur Lewis}, \omterm{def:b1:lewis-resonance-vsepr:formal-charge}{muatan formal}, dan resonansi
(\cref{ch:b1:lewis-resonance-vsepr}); $\mathrm{p}K_a$ dan kekuatan asam
dan basa (\cref{ch:b1:predominant-reaction}). Jilid sekolah (kelas 12)
menggambar \omterm{def:b1:electronic-effects:curly-arrow}{panah lengkung} dari \omterm{def:b1:electronic-effects:nucleophile}{nukleofil} ke \omterm{def:b1:electronic-effects:nucleophile}{elektrofil}.
\end{recall}

\section{Efek induktif}

\begin{definition}[Efek induktif]\label{def:b1:electronic-effects:inductive}
\emph{Efek induktif}\index{efek induktif} adalah pergeseran rapatan
elektron di sepanjang ikatan $\sigma$ ke arah atom atau gugus yang
elektronegatif. Gugus penarik (F, Cl, O, \ce{NO2}, \ce{CF3}) mempunyai
efek $-I$; gugus alkil dan atom bermuatan negatif mendorong elektron,
sebuah efek $+I$. Efek itu cepat melemah seiring bertambahnya banyaknya
ikatan di antaranya.
\end{definition}

\begin{proposition}[Efek induktif dan keasaman]\label{prop:b1:electronic-effects:inductive-pka}
Gugus penarik di dekat situs asam sebuah asam karboksilat menstabilkan
ion karboksilat dengan menyebarkan muatan negatifnya, dan menurunkan
$\mathrm{p}K_a$: makin banyak gugus penarik, makin dekat letaknya, makin
kuat asamnya.
\end{proposition}

\begin{proof}
$\mathrm{p}K_a$ mengukur posisi $\mathrm{RCOOH} \rightleftharpoons
\mathrm{RCOO^-} + \ce{H+}$. Gugus yang menarik rapatan elektron menjauh
dari karboksilat menyebarkan muatannya ke lebih banyak atom, yang
menurunkan energinya lebih banyak daripada energi asam netral;
kesetimbangannya bergeser ke kanan dan $K_a$ bertambah. Deret terukur
menegaskannya.
\end{proof}

\begin{omfigure}
\begin{tikzpicture}
\begin{axis}[
  xbar, width=0.78\linewidth, height=5.2cm, xmin=0, xmax=5.4,
  symbolic y coords={TFA,TCA,DCA,MCA,AA}, ytick=data,
  yticklabels={\ce{CF3COOH}, \ce{CCl3COOH}, \ce{CHCl2COOH}, \ce{CH2ClCOOH}, \ce{CH3COOH}},
  xlabel={$\mathrm{p}K_a$ pada \qty{25}{\celsius}}, bar width=9pt,
  nodes near coords, nodes near coords align={horizontal},
  every node near coord/.append style={font=\scriptsize},
  tick label style={font=\small}, label style={font=\small}, enlarge y limits=0.15]
\addplot[fill=omDef!45, draw=omDef] coordinates {(0.52,TFA) (0.51,TCA) (1.35,DCA) (2.87,MCA) (4.76,AA)};
\end{axis}
\end{tikzpicture}

{\small Asam-asam etanoat terhalogenasi. Setiap klor menurunkan
$\mathrm{p}K_a$; dengan tiga klor, asam itu sekitar sepuluh ribu kali
lebih kuat daripada asam etanoat. Asam trifluoro- dan trikloroetanoat
sama kuat dalam batas ketidakpastian pengukuran pada asam-asam yang
sekuat itu.}
% ledger: pka:ethanoic, pka:chloroethanoic, pka:dichloroethanoic, pka:trichloroethanoic, pka:trifluoroethanoic
\end{omfigure}

Gugus alkil bekerja ke arah sebaliknya, sedikit: asam propanoat
($\mathrm{p}K_a
= 4.89$) sedikit lebih lemah daripada asam etanoat (4.76), dan asam
metanoat (3.74), yang tidak mempunyai gugus alkil, lebih kuat daripada
keduanya.
% ledger: pka:propanoic, pka:methanoic

\section{Efek mesomeri}

Jika atom yang membawa \omterm{def:b1:lewis-resonance-vsepr:lewis}{pasangan elektron bebas}, atau muatan, terikat pada
sistem $\pi$, elektron-elektronnya dapat dipakai bersama dengan sistem
itu; struktur-struktur resonansi (\cref{ch:b1:lewis-resonance-vsepr})
memerikan pemakaian bersama itu.

\begin{definition}[Efek mesomeri, gugus donor dan gugus akseptor]\label{def:b1:electronic-effects:mesomeric}
\emph{Efek mesomeri}\index{efek mesomeri} adalah pergeseran rapatan
elektron melalui sistem $\pi$ terkonjugasi, yang diperlihatkan oleh
struktur-struktur resonansi. \emph{Gugus donor}\index{gugus donor}
memberikan \omterm{def:b1:lewis-resonance-vsepr:lewis}{sepasang elektron bebas} kepada sistem itu, sebuah efek $+M$
(\ce{-OH}, \ce{-O^-}, \ce{-NH2}, \ce{-OR}); \emph{gugus
akseptor}\index{gugus akseptor} menarik elektron dari sistem itu melalui
ikatan $\pi$ dengan atom elektronegatif, sebuah efek $-M$ (\ce{-NO2},
\ce{-C=O}, \ce{-C#N}).
\end{definition}

\begin{omfigure}
\setchemfig{atom sep=2em}
\schemestart
\chemfig{CH_3-C(=[1]O)-[7]O^{-}}
\arrow{<->}
\chemfig{CH_3-C(-[1]O^{-})=[7]O}
\schemestop
\hspace{2.5em}
\schemestart
\chemfig{*6(=-=(-O^{-})-=-)}
\arrow{<->}
\chemfig{*6((=O)-=-\chemabove{}{\scriptstyle\ominus}-=-)}
\schemestop

\medskip
\schemestart
\chemfig{H_3C-C(=[1]O)-[7]NH_2}
\arrow{<->}
\chemfig{H_3C-C(-[1]O^{-})=[7]NH_2^{+}}
\schemestop

{\small \omterm{def:b1:lewis-resonance-vsepr:resonance}{Struktur resonansi} ion etanoat (muatannya terbagi pada dua oksigen
yang ekuivalen), ion fenoksida (muatannya tersebar ke dalam cincin; satu
dari ketiga struktur cincin digambar), dan sebuah amida (\omterm{def:b1:lewis-resonance-vsepr:lewis}{pasangan
elektron bebas} nitrogen dipakai bersama dengan C=O, yang membuat amida
menjadi basa yang lemah).}
\end{omfigure}

\begin{proposition}[Efek mesomeri dan keasaman]\label{prop:b1:electronic-effects:mesomeric-pka}
Asam yang basa konjugasinya menyebarkan muatannya melalui resonansi lebih
kuat daripada asam yang basanya tidak dapat melakukannya: asam karboksilat
($\mathrm{p}K_a$ sekitar 4--5) jauh lebih kuat daripada alkohol (sekitar
16), dan fenol (10.0) terletak di antaranya. \omterm{def:b1:electronic-effects:mesomeric}{Gugus akseptor} pada posisi
\textit{orto} atau \textit{para} terhadap OH sebuah fenol menguatkannya
lebih banyak daripada pada posisi \textit{meta}.
\end{proposition}

\begin{proof}
Dalam alkoksida muatannya berada pada satu oksigen; dalam karboksilat
muatan itu terbagi pada dua oksigen; dalam fenoksida muatan itu terbagi
pada oksigen dan tiga karbon cincin, yang kurang elektronegatif daripada
oksigen, sebuah keuntungan yang lebih kecil. Gugus \ce{NO2} pada karbon
\textit{orto} atau \textit{para} dapat mengambil muatan itu sendiri
melalui \omterm{def:b1:lewis-resonance-vsepr:resonance}{struktur resonansi} dengan \ce{C=N+(-O^-)-O^-}; pada
\textit{meta}, tidak ada \omterm{def:b1:lewis-resonance-vsepr:resonance}{struktur resonansi} yang membawa muatan ke karbon
yang membawa \ce{NO2}, dan hanya \omterm{def:b1:electronic-effects:inductive}{efek induktif} yang tersisa. Nilai-nilai
terukurnya: 2-nitrofenol 7.23, 4-nitrofenol 7.16, 3-nitrofenol 8.36,
dibandingkan dengan 9.99 untuk fenol dan 15.9 untuk etanol.
\end{proof}
% ledger: pka:ethanol, pka:phenol, pka:2-nitrophenol, pka:3-nitrophenol, pka:4-nitrophenol

Penalaran yang sama berlaku untuk basa. Metilamina ($\mathrm{p}K_a$ ion
amoniumnya 10.66) adalah basa yang lebih kuat daripada amonia (9.25,
\cref{ch:b1:predominant-reaction}), karena gugus metil mendorong elektron
ke nitrogen; pada anilina \omterm{def:b1:lewis-resonance-vsepr:lewis}{pasangan elektron bebasnya} dipakai bersama
dengan cincin dan ion amoniumnya jauh lebih asam (4.60). Piridina (5.23)
mempertahankan \omterm{def:b1:lewis-resonance-vsepr:lewis}{pasangan elektron bebasnya}, tetapi pada nitrogen yang
terikat dalam cincin datar dengan sifat \textit{s} yang lebih besar,
sehingga kurang tersedia daripada pada amina.
% ledger: pka:methylammonium, pka:anilinium, pka:pyridinium, dfg:NH4+_ao

\section{Memutus ikatan: zat antara reaktif}

\begin{definition}[Homolisis dan heterolisis]\label{def:b1:electronic-effects:cleavage}
Sebuah \omterm{def:b1:lewis-resonance-vsepr:lewis}{ikatan kovalen} putus melalui \emph{homolisis}\index{homolisis}
jika setiap atom mempertahankan satu dari kedua elektron ikatan, yang
menghasilkan dua \omterm{def:b1:electronic-effects:intermediates}{radikal}, dan melalui
\emph{heterolisis}\index{heterolisis} jika satu atom mempertahankan
keduanya, yang menghasilkan sebuah kation dan sebuah anion.
\end{definition}

\begin{definition}[Zat antara reaktif]\label{def:b1:electronic-effects:intermediates}
\emph{Zat antara reaktif}\index{zat antara reaktif} adalah spesies yang
terbentuk pada satu tahap sebuah mekanisme dan terpakai pada tahap
sesudahnya, yang terlalu reaktif untuk menumpuk
(\cref{ch:b1:elementary-steps}). \emph{Karbokation}\index{karbokation}
mempunyai karbon dengan tiga ikatan dan muatan positif (enam \omterm{def:b1:quantum-numbers:configuration}{elektron
valensi}); \emph{karbanion}\index{karbanion} mempunyai karbon dengan tiga
ikatan, \omterm{def:b1:lewis-resonance-vsepr:lewis}{sepasang elektron bebas}, dan muatan negatif;
\emph{radikal}\index{radikal} mempunyai atom dengan \omterm{def:b1:quantum-numbers:unpaired}{elektron tak
berpasangan}.
\end{definition}

\begin{definition}[Kelas sebuah karbon]\label{def:b1:electronic-effects:carbon-class}
Karbon yang terikat pada satu karbon lain adalah \emph{karbon
primer}\index{karbon primer}, pada dua karbon lain \emph{karbon
sekunder}\index{karbon sekunder}, pada tiga karbon lain \emph{karbon
tersier}\index{karbon tersier}. \omterm{def:b1:electronic-effects:intermediates}{Karbokation}, \omterm{def:b1:electronic-effects:intermediates}{radikal}, atau haloalkana
mengambil kelas karbon yang bersangkutan.
\end{definition}

\begin{omfigure}
\begin{tikzpicture}[font=\small]
  % carbocation: planar, empty p orbital
  \begin{scope}
    \fill[omProp!15, draw=omProp] (0,0.15) ellipse (0.18 and 0.75);
    \fill[omProp!15, draw=omProp] (0,-0.15) ellipse (0.18 and 0.75);
    \draw[thick] (0,0) -- (-1.0,-0.25) node[left] {R};
    \draw[thick] (0,0) -- (0.95,-0.35) node[right] {R};
    \draw[thick] (0,0) -- (0.35,0.45) node[above right] {R};
    \node[fill=white, inner sep=1pt] at (0,0) {\ce{C+}};
    \node at (0,-1.75) {karbokation: datar,};
    \node at (0,-2.15) {orbital \textit{p} kosong};
  \end{scope}
  % radical: nearly planar, one electron
  \begin{scope}[xshift=4.6cm]
    \draw[omMeth] (0,0.15) ellipse (0.18 and 0.75);
    \draw[omMeth] (0,-0.15) ellipse (0.18 and 0.75);
    \fill (0,0.55) circle (0.05);
    \draw[thick] (0,0) -- (-1.0,-0.25) node[left] {R};
    \draw[thick] (0,0) -- (0.95,-0.35) node[right] {R};
    \draw[thick] (0,0) -- (0.35,0.45) node[above right] {R};
    \node[fill=white, inner sep=1pt] at (0,0) {C};
    \node at (0,-1.75) {radikal: hampir datar,};
    \node at (0,-2.15) {satu elektron};
  \end{scope}
  % carbanion: pyramidal with a lone pair
  \begin{scope}[xshift=9.2cm]
    \draw[thick] (0,0) -- (-0.95,-0.55) node[left] {R};
    \draw[thick] (0,0) -- (0.95,-0.55) node[right] {R};
    \draw[thick] (0,0) -- (0.2,-0.9) node[below] {R};
    \fill[omDef!15, draw=omDef] (0,0.55) ellipse (0.22 and 0.45);
    \fill (-0.07,0.65) circle (0.045); \fill (0.07,0.65) circle (0.045);
    \node[fill=white, inner sep=1pt] at (0,0) {\ce{C-}};
    \node at (0,-1.75) {karbanion: piramida,};
    \node at (0,-2.15) {pasangan elektron bebas};
  \end{scope}
\end{tikzpicture}

{\small Geometri ketiga \omterm{def:b1:electronic-effects:intermediates}{zat antara reaktif} karbon (R: H atau gugus alkil).
\omterm{def:b1:electronic-effects:intermediates}{Karbokation} yang datar dapat diserang dari kedua sisinya, hal yang
penting bagi stereokimia pada
\cref{ch:b1:nucleophilic-substitution}.}
\end{omfigure}

\begin{proposition}[Kestabilan karbokation]\label{prop:b1:electronic-effects:carbocation-order}
\omterm{def:b1:electronic-effects:intermediates}{Karbokation} distabilkan oleh gugus-gugus yang memberinya elektron:
tersier $>$ sekunder $>$ primer $>$ metil; \omterm{def:b1:electronic-effects:intermediates}{karbokation} di sebelah ikatan
rangkap dua (alilik) atau cincin benzena (benzilik) distabilkan lebih
lanjut melalui resonansi, dan \omterm{def:b1:electronic-effects:intermediates}{karbokation} di sebelah atom yang membawa
\omterm{def:b1:lewis-resonance-vsepr:lewis}{pasangan elektron bebas} (\ce{R-O-CH2+}) lebih lagi.
\end{proposition}

\begin{proof}
Setiap gugus alkil mendorong rapatan elektron ke arah karbon yang miskin
elektron (efek $+I$, ditambah interaksi terkait antara ikatan-ikatan
C--H-nya dan orbital kosong itu, yang dipaparkan dalam jilid Tahun~2).
Kation alilik mempunyai dua \omterm{def:b1:lewis-resonance-vsepr:resonance}{struktur resonansi} yang menempatkan muatan
pada salah satu karbon ujung; kation benzilik mempunyai empat, tiga di
dalam cincin; \omterm{def:b1:lewis-resonance-vsepr:lewis}{sepasang elektron bebas} oksigen membentuk ikatan keempat
pada karbon kationik, \ce{R-O+=CH2}, yang setiap atomnya mempunyai oktet.
\end{proof}

\omterm{def:b1:electronic-effects:intermediates}{Karbanion} mengikuti urutan yang berlawanan (metil $>$ primer $>$ sekunder
$>$ tersier) dan distabilkan oleh \omterm{def:b1:electronic-effects:mesomeric}{gugus akseptor}; \omterm{def:b1:electronic-effects:intermediates}{radikal} mengikuti
urutan yang sama dengan \omterm{def:b1:electronic-effects:intermediates}{karbokation}, dengan perbedaan yang kurang
mencolok.

\section{Panah lengkung, nukleofil, dan elektrofil}

\begin{definition}[Panah lengkung]\label{def:b1:electronic-effects:curly-arrow}
Dalam sebuah mekanisme, \emph{panah lengkung}\index{panah lengkung}
memperlihatkan perpindahan sepasang elektron: panah itu berawal pada
\omterm{def:b1:lewis-resonance-vsepr:lewis}{pasangan elektron bebas} atau ikatan dan berakhir pada sebuah atom
(tempat terbentuk \omterm{def:b1:lewis-resonance-vsepr:lewis}{pasangan elektron bebas} atau ikatan baru) atau di
antara dua atom (sebuah ikatan baru). \emph{Panah berkepala
setengah}\index{panah berkepala setengah} memperlihatkan perpindahan satu
elektron, dalam reaksi \omterm{def:b1:electronic-effects:intermediates}{radikal}.
\end{definition}

\begin{method}[Menggambar panah lengkung]\label{met:b1:electronic-effects:curly-arrows}
\begin{enumerate}
  \item Mulailah dari elektron: \omterm{def:b1:lewis-resonance-vsepr:lewis}{pasangan elektron bebas}, ikatan $\pi$,
        atau ikatan $\sigma$, jangan pernah dari muatan positif atau atom
        tanpa elektron.
  \item Akhiri di tempat pasangan itu pergi: ikatan baru antara kedua
        atom, atau \omterm{def:b1:lewis-resonance-vsepr:lewis}{pasangan elektron bebas} pada atom yang lebih
        elektronegatif ketika sebuah ikatan putus.
  \item Hitunglah muatan sesudah setiap tahap: atom yang memberikan
        \omterm{def:b1:lewis-resonance-vsepr:lewis}{pasangan elektron bebas} kepada ikatan baru memperoleh $+1$, atom
        yang menerima pasangan sebagai \omterm{def:b1:lewis-resonance-vsepr:lewis}{pasangan elektron bebas} memperoleh
        $-1$; muatan totalnya kekal.
  \item Jangan pernah melampaui oktet pada C, N, O: ketika sepasang
        elektron datang, pasangan lain harus pergi.
\end{enumerate}
\end{method}

\begin{omfigure}
\setchemfig{atom sep=2.2em}
\schemestart
\chemfig{H_3C-C(-[2]H)(-[6]H)-[@{b}]@{br}Br}
\arrow{->[heterolisis]}[,1.5]
\chemfig{H_3C-C(-[2]H)(-[6]H)^{+}}\+\chemfig{Br^{-}}
\schemestop
\chemmove{\draw[->, shorten <=2pt, shorten >=2pt](b) .. controls +(90:6mm) and +(90:6mm) .. (br);}

\medskip
\schemestart
\chemfig{H@{o}O^{-}}\+\chemfig{@{h}H-@{oc}O-H}
\arrow{<=>}
\chemfig{H_2O}\+\chemfig{HO^{-}}
\schemestop
\chemmove{\draw[->, thick, shorten <=2pt, shorten >=2pt](o) .. controls +(-80:8mm) and +(-100:8mm) .. (h);}
\par\vspace{5mm}

{\small \omterm{def:b1:electronic-effects:curly-arrow}{Panah lengkung}. Atas: \omterm{def:b1:electronic-effects:cleavage}{heterolisis} ikatan C--Br, dengan pasangan
elektron menuju brom. Bawah: perpindahan proton dari air ke hidroksida,
yang ditulis sebagai serangan \omterm{def:b1:lewis-resonance-vsepr:lewis}{pasangan elektron bebas} basa pada hidrogen
(ikatan O--H air lalu putus, dengan pasangannya tetap pada oksigen;
panah kedua ditinggalkan untuk \cref{exo:b1:electronic-effects:4}).}
\end{omfigure}

\begin{definition}[Nukleofil, elektrofil, gugus pergi]\label{def:b1:electronic-effects:nucleophile}
\emph{Nukleofil}\index{nukleofil} adalah spesies yang memberikan sepasang
elektron kepada atom selain hidrogen untuk membentuk ikatan (sebuah \omterm{def:b1:lewis-resonance-vsepr:lewis-acid}{basa
Lewis}); \emph{elektrofil}\index{elektrofil} adalah spesies yang menerima
pasangan semacam itu (sebuah \omterm{def:b1:lewis-resonance-vsepr:lewis-acid}{asam Lewis}, sering kali karbon yang membawa
muatan parsial positif). \emph{Nukleofilisitas}\index{nukleofilisitas}
sebuah spesies adalah kereaktifannya sebagai nukleofil, yang diukur
dengan \omterm{def:b1:rate-laws:order}{tetapan laju}. \emph{Gugus pergi}\index{gugus pergi} adalah gugus
yang pergi bersama \omterm{def:b1:lewis-resonance-vsepr:lewis}{pasangan elektron ikatan} ketika sebuah ikatan dengan
karbon putus secara heterolitik.
\end{definition}

\begin{proposition}[Nukleofilisitas dan kebasaan]\label{prop:b1:electronic-effects:nucleophilicity-basicity}
Untuk \omterm{def:b1:electronic-effects:nucleophile}{nukleofil-nukleofil} yang menyerang melalui atom yang sama,
\omterm{def:b1:electronic-effects:nucleophile}{nukleofilisitas} mengikuti kebasaan: \ce{CH3O-} $>$ \ce{OH-} $>$
\ce{CH3COO-} $>$ \ce{H2O}. Hal ini tidak berlaku sepanjang satu kolom
tabel periodik dalam \omterm{def:b1:intermolecular-forces-solvents:solvent-classes}{pelarut protik}, tempat \omterm{def:b1:electronic-effects:nucleophile}{nukleofil} yang besar, mudah
terpolarisasi, dan lemah tersolvasi bereaksi lebih cepat walaupun
merupakan basa yang lebih lemah (\ce{I-} $>$ \ce{Br-} $>$ \ce{Cl-} $>$
\ce{F-}), dan untuk basa yang meruah, yang merupakan \omterm{def:b1:electronic-effects:nucleophile}{nukleofil} yang
buruk.
\end{proposition}

\begin{proof}
Kebasaan mengukur kesetimbangan serangan pada proton, \omterm{def:b1:electronic-effects:nucleophile}{nukleofilisitas}
mengukur laju serangan pada karbon; keduanya bertambah dengan
ketersediaan pasangan elektron, itulah sebabnya keduanya sejalan untuk
satu atom. Lajunya juga bergantung pada energi yang diperlukan untuk
melepaskan pelarut dari \omterm{def:b1:electronic-effects:nucleophile}{nukleofil} dan pada distorsi awan elektronnya ke
arah karbon yang padat: anion kecil seperti \ce{F-} terikat kuat pada
molekul-molekul \omterm{def:b1:intermolecular-forces-solvents:solvent-classes}{pelarut protik}
(\cref{ch:b1:intermolecular-forces-solvents}), dan gugus-gugus yang
meruah menghalangi pendekatan ke karbon tetapi tidak ke proton.
\end{proof}

\omterm{def:b1:electronic-effects:nucleophile}{Gugus pergi} yang baik adalah \omterm{def:b1:predominant-reaction:strong-acid}{basa lemah}, yang stabil setelah membawa
pasangan elektron itu: \ce{I-}, \ce{Br-}, \ce{Cl-}, air, sulfonat;
\ce{OH-} dan \ce{RO-}, basa-basa kuat, adalah \omterm{def:b1:electronic-effects:nucleophile}{gugus pergi} yang buruk ---
itulah sebabnya alkohol harus diaktifkan sebelum substitusi
(\cref{ch:b1:alcohol-activation}).

\section{Latihan}

\begin{exercise}[$\star$]\label{exo:b1:electronic-effects:1}
Golongkan setiap karbon 2-metilbutana dan 2-bromo-2-metilpropana sebagai
primer, sekunder, tersier, atau kuaterner (terikat pada empat karbon).
\omterm{def:b1:electronic-effects:intermediates}{Karbokation} manakah yang paling mudah terbentuk dari setiap bromida
berumus \ce{C4H9Br}?
\end{exercise}

\begin{exercise}[$\star$]\label{exo:b1:electronic-effects:2}
Dari nilai-nilai $\mathrm{p}K_a$, hitunglah rasio \omterm{def:b1:predominant-reaction:acidity-constant}{tetapan keasaman} asam
kloroetanoat dan asam etanoat, serta asam trifluoroetanoat dan asam
etanoat.
\end{exercise}

\begin{exercise}[$\star$]\label{exo:b1:electronic-effects:3}
Gambarlah struktur-struktur resonansi ion etanoat, ion 4-nitrofenoksida
(dengan muatan pada salah satu oksigen gugus nitro), dan kation alil
\ce{CH2=CH-CH2+}.
\end{exercise}

\begin{exercise}[$\star$]\label{exo:b1:electronic-effects:4}
Lengkapilah panah-panah lengkung perpindahan proton pada gambar itu dan
pada \ce{NH3 + H3O+ -> NH4+ + H2O}. Periksalah muatan-muatannya.
\end{exercise}

\begin{exercise}[$\star\star$]\label{exo:b1:electronic-effects:5}
Urutkan menurut keasaman yang bertambah: etanol, fenol, asam etanoat,
4-nitrofenol, asam trikloroetanoat, air ($\mathrm{p}K_a$ pasangan
\ce{H2O}/\ce{OH-}, 14.0). Berikan alasan setiap langkah dengan argumen
induktif atau mesomeri.
\end{exercise}

\begin{exercise}[$\star\star$]\label{exo:b1:electronic-effects:6}
Urutkan menurut kebasaan yang bertambah: anilina, metilamina, amonia,
piridina. Jelaskan letak anilina dengan sebuah \omterm{def:b1:lewis-resonance-vsepr:resonance}{struktur resonansi}.
\end{exercise}

\begin{exercise}[$\star\star$]\label{exo:b1:electronic-effects:7}
Mengapa 3-nitrofenol lebih lemah daripada 4-nitrofenol tetapi tetap lebih
kuat daripada fenol? Gambarlah struktur-struktur resonansi yang mendukung
kedua pernyataan itu.
\end{exercise}

\begin{exercise}[$\star\star$]\label{exo:b1:electronic-effects:8}
Urutkan \omterm{def:b1:electronic-effects:intermediates}{karbokation} \ce{CH3+}, \ce{(CH3)3C+}, \ce{CH2=CH-CH2+},
\ce{CH3CH2+}, \ce{C6H5CH2+}, \ce{CH3OCH2+}, lalu berikan alasannya.
\end{exercise}

\begin{exercise}[$\star\star$]\label{exo:b1:electronic-effects:9}
Tentukan \omterm{def:b1:electronic-effects:nucleophile}{nukleofil} dan \omterm{def:b1:electronic-effects:nucleophile}{elektrofil} dalam: \ce{CH3Br + OH-};
\ce{H2O + H+}; etanal dengan ion sianida; \ce{BF3 + NH3}. Gambarlah \omterm{def:b1:electronic-effects:curly-arrow}{panah
lengkungnya}.
\end{exercise}

\begin{exercise}[$\star\star\star$]\label{exo:b1:electronic-effects:10}
Sebuah larutan mengandung \qty{0.010}{mol/L} asam trikloroetanoat.
Apakah rumus \omterm{def:b1:predominant-reaction:strong-acid}{asam lemah} $\mathrm{pH} = \frac12(\mathrm{p}K_a - \log C)$
berlaku? Hitunglah pH-nya dengan benar dan \omterm{def:b1:predominant-reaction:dissociation}{derajat disosiasinya}.
\end{exercise}

\begin{exercise}[$\star\star\star$]\label{exo:b1:electronic-effects:11}
Nilai $\mathrm{p}K_a$ alkohol yang diukur di dalam air bertambah dari
metanol (15.5) ke etanol (15.9), propan-2-ol (17.1), dan
2-metilpropan-2-ol (19.2). Berikan dua penjelasan, satu melalui \omterm{def:b1:electronic-effects:inductive}{efek
induktif} gugus-gugus alkil pada alkoksida, satu melalui \omterm{def:b1:intermolecular-forces-solvents:dissolution}{solvasi}
alkoksida.
% ledger: pka:methanol, pka:ethanol, pka:2-propanol, pka:tBuOH
\end{exercise}

\begin{exercise}[$\star\star\star$]\label{exo:b1:electronic-effects:12}
Jelaskan mengapa amida tidak bersifat basa pada nitrogennya dan tidak
nukleofilik, sedangkan amina bersifat keduanya, dengan memakai \omterm{def:b1:lewis-resonance-vsepr:resonance}{struktur
resonansi} amida. Atom manakah pada amida yang terprotonasi dalam \omterm{def:b1:predominant-reaction:strong-acid}{asam
kuat}?
\end{exercise}

\section{Soal: Mengapa Asam Trifluoroetanoat Begitu Kuat}

\begin{problem}[{Soal akhir pekan --- deret induktif asam-asam
terhalogenasi, karboksilat dibandingkan dengan alkoksida, ketiga
nitrofenol, sebuah perhitungan reaksi predominan, dan rasio tetapan
keasaman asam trifluoroetanoat dan asam etanoat}]
\label{pb:b1:electronic-effects:1}
Data ($\mathrm{p}K_a$ pada \qty{25}{\celsius}): asam etanoat 4.76,
kloroetanoat 2.87, dikloroetanoat 1.35, trikloroetanoat 0.51,
trifluoroetanoat 0.52, etanol 15.9, fenol 9.99, 2-nitrofenol 7.23,
3-nitrofenol 8.36, 4-nitrofenol 7.16; $\mathrm{p}K_e = 14.00$.

\medskip
\noindent\textbf{Bagian I --- Deret induktif.}
\begin{enumerate}
  \item Hitunglah rasio $K_a$ untuk setiap langkah dari asam etanoat
        sampai asam trikloroetanoat.
  \item Apakah pengaruh setiap klor sama? Berikan komentar.
  \item Jelaskan kecenderungan itu dengan \omterm{def:b1:electronic-effects:inductive}{efek induktif}.
  \item Mengapa klor pada karbon sebuah rantai yang lebih panjang yang
        letaknya paling jauh dari COOH hampir tidak berpengaruh?
  \item Fluorin lebih elektronegatif daripada klor. Apa yang ditunjukkan
        perbandingan asam trifluoro- dan trikloro-, dan mengapa
        $\mathrm{p}K_a$ semacam itu sulit diukur dengan teliti?
  \item Dalam bentuk apakah setiap asam itu berada pada pH 3?
\end{enumerate}

\medskip
\noindent\textbf{Bagian II --- Karboksilat dan alkoksida.}
\begin{enumerate}[resume]
  \item Gambarlah kedua \omterm{def:b1:lewis-resonance-vsepr:resonance}{struktur resonansi} ion etanoat.
  \item Apa yang diramalkan keduanya tentang panjang kedua ikatan C--O
        ion itu?
  \item Mengapa ion etoksida tidak mempunyai struktur-struktur yang
        ekuivalen?
  \item Hitunglah rasio $K_a$ asam etanoat dan etanol.
  \item Tempatkan fenol di antara keduanya dan berikan alasan dengan
        struktur-struktur resonansi fenoksida.
  \item Manakah di antara etanol, fenol, dan asam etanoat yang bereaksi
        hampir tuntas dengan larutan hidrogenkarbonat ($\mathrm{p}K_a$ 6.37
        untuk \ce{CO2,H2O}/\ce{HCO3-})?
\end{enumerate}

\medskip
\noindent\textbf{Bagian III --- Nitrofenol.}
\begin{enumerate}[resume]
  \item Urutkan ketiga nitrofenol dan fenol menurut keasamannya.
  \item Gambarlah sebuah \omterm{def:b1:lewis-resonance-vsepr:resonance}{struktur resonansi} 4-nitrofenoksida dengan
        muatan pada gugus nitro.
  \item Tunjukkan bahwa struktur semacam itu tidak ada untuk
        3-nitrofenoksida.
  \item Mengapa 3-nitrofenol tetap lebih asam daripada fenol?
  \item 2-Nitrofenol sedikit lebih lemah daripada 4-nitrofenol walaupun
        keduanya mempunyai \omterm{def:b1:electronic-effects:mesomeric}{efek mesomeri}; usulkan sebuah penjelasan yang
        melibatkan OH-nya dan gugus nitro di dekatnya.
  \item Pada pH 7.5, nitrofenol manakah yang sebagian besar berada dalam
        bentuk anion? Mengapa hal itu menjadikan 4-nitrofenol sebuah
        \omterm{def:b1:titration-methods:indicator}{indikator warna}?
\end{enumerate}

\medskip
\noindent\textbf{Bagian IV --- Asam trifluoroetanoat di dalam air.}
\begin{enumerate}[resume]
  \item Hitunglah pH asam trifluoroetanoat \qty{0.010}{mol/L} (jangan
        menganggap disosiasinya lemah).
  \item Hitunglah pH asam etanoat \qty{0.010}{mol/L}.
  \item Tuliskan reaksi asam trifluoroetanoat dengan natrium etanoat dan
        hitunglah tetapannya.
  \item Nyatakan rasio $K_a(\ce{CF3COOH})/K_a(\ce{CH3COOH})$, sebagai
        pangkat sepuluh dan sebagai bilangan.
\end{enumerate}
\end{problem}
